\chapter{Solid State Theory: Midterm, 29 April 2016}
\label{q:theory-2016-04-midterm}
\begin{sourceNote}
The supplied printed sheet contains all twelve questions, whose weights total 100\%. Mathematical subscripts and line breaks have been typeset consistently; handwritten answers are excluded.
\end{sourceNote}

\section*{1. (15\%) True or false. Explain for your answer.}
\begin{enumerate}[label=(\alph*)]
\item (5\%) The determinantal wavefunction with spin-orbitals satisfying the Hartree--Fock equations is the exact solution of the Schrödinger equation for many interacting electrons.
\item (5\%) Consider a homogeneous electron gas embedded in a uniform positive charge background neutralizing the negative charge of the electron gas. The electron gas tends to be fully spin-polarized as the electron density increases.
\item (5\%) A crystal with an odd number of electrons in each primitive cell cannot be an electrical insulator.
\end{enumerate}

\section*{2. (5\%) Hamiltonian}
Express the Hamiltonian for a system of $N$ interacting electrons in the presence of nuclei fixed in a configuration. Explain each term in your expression.

\section*{3. (10\%) Trial wave functions}
What are the trial wave functions for the many-electron system in Problem 2 in the Hartree theory and in the Hartree--Fock theory, respectively?

\section*{4. (10\%) Hartree--Fock equations}
Express the Hartree--Fock equations and explain each term in your expression. Roughly describe how the equations result from the trial wave functions.

\section*{5. (5\%) Hohenberg--Kohn functional}
Express and explain the Hohenberg--Kohn functional.

\section*{6. (10\%) Kohn--Sham equations}
Express the Kohn--Sham equations and explain each term in your expression. Roughly describe how the equations result from the Hohenberg--Kohn functional.

\section*{7. (5\%) Variational collapse}
Describe and explain the variational collapse in the expansion of crystal states with pure plane waves.

\section*{8. (5\%) Modified plane waves}
How to modify the basis functions of plane waves to cure the variational collapse problem?

\section*{9. (5\%) Pseudo wavefunction}
Express the relation between a crystal state wavefunction $\psi_{\vct{k}}$ and its pseudo wavefunction $\varphi_{\vct{k}}$.

\section*{10. (10\%) Effective Hamiltonian}
The Schrödinger equation for a crystal state $\psi_{\vct{k}}$ can become an effective Schrödinger equation for the pseudo wavefunction $\varphi_{\vct{k}}$ with an effective Hamiltonian $H^{(\mathrm{eff})}$.
\begin{enumerate}[label=(\alph*)]
\item (5\%) Express and explain the effective Hamiltonian $H^{(\mathrm{eff})}$.
\item (5\%) Roughly describe how to solve the effective Schrödinger equation using the basis of plane waves.
\end{enumerate}

\section*{11. (5\%) Sodium valence energy}
Conventionally, the $3s$ valence electron of the sodium atom is regarded to experience an effective Coulomb potential of a point-like charge, as the $1s$ electron of the hydrogen atom. Explain the fact that the $3s$ orbital energy of the sodium atom is about $-5.14\,\mathrm{eV}$ instead of $-13.6\,\mathrm{eV}$ for the $1s$ orbital energy of the hydrogen atom.

\section*{12. (5\%) Pseudo atom}
Explain the term ``pseudo atom.''
